\documentclass[11pt,a4paper]{article}

\usepackage[margin=1in]{geometry}
\usepackage{amsmath,amssymb}
\usepackage{booktabs}
\usepackage{hyperref}
\usepackage{microtype}
\usepackage{enumitem}
\usepackage{natbib}
\usepackage{xcolor}

\hypersetup{
  colorlinks=true,
  linkcolor=black,
  citecolor=black,
  urlcolor=blue!50!black
}

\title{GPNEC: A Metal State Automaton on Apple Silicon\\
\large Verified Lattice Boltzmann Parity, Dual-Geometry Routing\\
under Matched Policies, and Honest Performance Framing}
\author{Digital Defiance}
\date{August 2026}

\begin{document}
\maketitle

\begin{abstract}
We present \textbf{GPNEC}, an open Apple Silicon engine for hand-crafted
state automata $X_{t+1}=\Phi(X_t,W,\theta)$ using Metal compute kernels and,
where topology is a dense linear map, MPSGraph. This report focuses on what
we can verify with reproducible CLI gates on an Apple~M4~Max.

For D2Q9 Lattice Boltzmann (BGK), a single-thread CPU reference that mirrors
the Metal collide--stream kernels matches GPU state to relative
$\mathrm{L}_2\sim3\times10^{-7}$. We report Metal throughput in MLups and
estimate effective memory traffic, and we explicitly treat the single-thread
CPU mirror as an \emph{accuracy gold standard}, not as a competitive CFD
baseline. Headline ``speedup vs one CPU thread'' figures are therefore
secondary and carefully scoped.

For dual-geometry greedy routing, Metal and CPU hop updates agree with zero
field mismatches under lockstep comparison. Under a \textbf{symmetric}
post-crash protocol---identical silent-retry and respawn on both
metrics---Poincar\'e cumulative deliveries remain higher than Euclidean on
our dual embedding (about $1.8\times$ in the primary table). We also report
the Route app's asymmetric ``sandbox'' protocol for completeness, and we
state plainly that it handicaps Euclidean traffic and must not be used as an
embedding-only resilience claim.

Source, \texttt{gpnec verify}, and paired \texttt{route-verify --policy both}
gates are released with the measurements archived under
\texttt{papers/gpnec/}.
\end{abstract}

\tableofcontents
\newpage

\section{Introduction}
\label{sec:intro}

GPNEC evaluates local update rules on Apple Silicon without an ML training
stack: no Core~ML, MLX, or PyTorch in the core loop. Geometry lives in $W$,
transition physics or routing policy in a Metal kernel $\Phi$, and the host
API stays stable across adapters.

This revision addresses methodological gaps that would (correctly) draw
reviewer criticism in a systems or computational-physics venue:
\begin{enumerate}[leftmargin=*]
  \item Routing comparisons now lead with a \textbf{symmetric} crash
        control (identical retry/respawn). The UI ``freeze Euclidean''
        protocol is demoted to a demo narrative.
  \item LBM performance is framed primarily as Metal MLups and bandwidth
        utilization against a kernel-faithful CPU accuracy reference---not
        as a $2000\times$ claim against a competitive multi-core CFD code.
  \item We separate MPSGraph topology paths from custom Metal LBM/route
        kernels, document Poincar\'e float32 clamping, and state LBM Mach /
        boundary assumptions for the verified bench path.
\end{enumerate}

\section{System architecture}
\label{sec:system}

\subsection{Automaton API}
\begin{equation}
  X_{t+1} = \Phi(X_t, W, \theta).
\end{equation}
State lives in double-buffered \texttt{MTLBuffer}s on unified memory.

\subsection{Where MPSGraph ends and Metal begins}
The marketing phrase ``Metal/MPSGraph automaton'' covers two execution
paths that must not be conflated:

\begin{center}
\begin{tabular}{@{}llp{0.48\textwidth}@{}}
\toprule
Path & Mechanism & Used by \\
\midrule
Topology matmul & \texttt{MPSGraph} batched $Y\!=\!WX$ then custom $\Phi$
  & Subspace / hyperbolic CLI adapters; optional dense LBM cost-model check \\
Collide+stream & Two custom Metal kernels (\texttt{lbm\_bgk\_collide},
  \texttt{lbm\_stream\_bounce} in \texttt{PhiLBM.metal}); MPSGraph is
  \emph{not} in the hot loop
  & Production LBM bench, \texttt{verify-lbm}, Fluid app \\
Dual greedy route & Custom Metal kernel \texttt{dual\_greedy\_route\_step}
  & Route app, \texttt{verify-route}, \texttt{route-verify} \\
\bottomrule
\end{tabular}
\end{center}

In short: MPSGraph manages dense linear topology where that is the natural
$W$; the headline LBM and dual-router results in this paper are custom Metal
kernels.

\section{Lattice Boltzmann}
\label{sec:lbm}

\subsection{Model, Mach number, and boundaries}
We use D2Q9 BGK~\cite{succi2001,kruger2017} with lattice sound speed
$c_s=1/\sqrt{3}$. Default demo inlet $u_x=0.08$ gives
\begin{equation}
  \mathrm{Ma} \approx \frac{u_x}{c_s} \approx 0.14,
\end{equation}
comfortably in the low-Mach regime where BGK is commonly applied. Relaxation
time $\tau=0.56$ (clamped $\ge0.51$ in kernels).

The stream kernel applies: periodic boundaries in $y$; approximate bounce on
open $x$ faces; mid-grid bounce-back from solid masks; and a left-column
velocity inlet that overwrites populations with equilibrium at
$(\rho,u)=(1,(u_x,0))$. Dye (channel~9) is advected upwind. The
\textbf{accuracy/bench path} used in tables below disables solids and uses
matched inlet settings so Metal and CPU see identical IC/BC; the Fluid UI
adds cylinder solids for visualization without changing the verify
contract.

\subsection{CPU reference (accuracy gold, not a CFD champion)}
\texttt{cpu-d2q9} is a single-thread Swift implementation that mirrors
\texttt{PhiLBM.metal} channel layout and arithmetic. Its job is
\textbf{algorithmic parity} for relative $\mathrm{L}_2$ tests. It is
intentionally not NEON-tuned multipole CFD. Optional \texttt{cpu-mt}
parallelizes over rows; on this dense stencil it is typically
\emph{slower} than single-thread (memory/contention), so we report it only
to discourage using a weak multi-thread straw man as a ``fixed'' baseline.

We therefore \textbf{do not} treat ``Metal / single-thread CPU'' ratios as
the primary systems claim. Primary claims are: (i)~Metal$\equiv$CPU state
parity, and (ii)~achieved Metal MLups with an order-of-magnitude bandwidth
estimate.

\subsection{Accuracy protocol}
\texttt{gpnec verify-lbm} compares full 10-channel state after $N$ steps.
Gate: relative $\mathrm{L}_2\le10^{-4}$ (observed $\sim10^{-7}$), finite
values, and rest-equilibrium drift on CPU.

\subsection{Throughput and bandwidth estimate}
Let $N_c=W\cdot H$ cells and $C=10$ float32 channels (D2Q9+dye). One
collide+stream step performs two full read/write passes over the state in
our ping-pong implementation, i.e.\ about
\begin{equation}
  T \approx 4\cdot N_c\cdot C\cdot 4\ \mathrm{bytes}
  = 16\,N_c\,C\ \mathrm{bytes}
\end{equation}
of explicit buffer traffic (ignoring solids texture and incidental reads).
For $256^2$ this is $\approx10.5$~MiB/step. At $69.9\,\mu\mathrm{s}$/step
(Table~\ref{tab:lbm-thru}),
\begin{equation}
  B_{\mathrm{eff}} \approx \frac{10.5\times2^{20}}{69.9\times10^{-6}}
  \approx 150\ \mathrm{GB/s}.
\end{equation}
Published M4~Max unified memory bandwidth is on the order of
$500{+}\,\mathrm{GB/s}$; $\sim150\,\mathrm{GB/s}$ is therefore only a
modest fraction of peak. At these lattice sizes the working set
($\sim2.5$~MiB) fits in caches, so the step time is plausibly limited by
\textbf{kernel launch / command-buffer overhead and latency} rather than
sustained DRAM bandwidth. Larger lattices, fused collide--stream, or
persistent command buffers would be needed before claiming a bandwidth
roofline result. We report this estimate so readers do not over-interpret
absolute MLups.

\section{Dual-geometry routing}
\label{sec:route}

\subsection{Dual embedding}
Sample Poincar\'e-disk points; build symmetrized $K$-NN edges in Poincar\'e
distance; embed the \emph{same} undirected graph in the plane via landmark
MDS; select a crash set of size $C$ along a vertical mid-strip of the
Euclidean layout (Route app backbone cut), informed by sampled betweenness.

\subsection{Greedy hop and float32 Poincar\'e clamping}
One Metal thread per packet forwards to the live neighbor minimizing
distance to the destination. Poincar\'e distance matches the CPU reference:
\begin{align}
  &\text{if }\lVert a\rVert\ge0.999,\quad a\leftarrow a\cdot0.999/\lVert a\rVert
    \quad\text{(likewise for $b$)},\\
  &d=\log\!\bigl(\alpha+\sqrt{\max(\alpha^2-1,0)}\bigr),\quad
    \alpha=1+\frac{2\lVert a-b\rVert^2}{\max(1-\lVert a\rVert^2,10^{-6})\,
    \max(1-\lVert b\rVert^2,10^{-6})}.
\end{align}
Clamping keeps points strictly inside the open disk so float32
denominators and $\mathrm{acosh}$-equivalent logs remain finite near the
boundary. This is a reproducibility detail, not a claim of exact hyperbolic
isometry under all float32 roundings.

\subsection{Metal$\equiv$CPU lockstep}
With \texttt{stuckRetryHops}=0, Metal and CPU steppers are compared field-wise
each tick. Zero mismatches is the correctness gate for the hop kernel.

\subsection{Crash protocols: symmetric (primary) vs sandbox (narrative)}
\label{sec:crash-protocols}

\paragraph{Symmetric (publication control).}
After the crash, \emph{both} metrics keep \texttt{stuckRetryHops}=40 and
respawn finished slots. Any delivery gap is attributable to geometry and
greedy progress on the damaged graph under matched mechanics.

\paragraph{Sandbox (UI narrative only).}
Matches the Route app: Euclidean sets \texttt{stuckRetryHops}=0 and stops
respawning; Poincar\'e keeps retry+respawn. This \textbf{handicaps}
Euclidean traffic by construction. It is useful for a visual demo and is
reported for fidelity to the shipped UI, but it is \textbf{not} evidence
that the embedding alone explains large delivery multiples.

\texttt{gpnec verify-route} requires the symmetric gate for
\texttt{verified: true}; sandbox results are printed alongside.
\texttt{route-verify --policy both} prints both reports.

\section{Experimental setup}
\label{sec:setup}

Host: Apple~M4~Max, macOS, Swift~6 / Metal. Transcripts:
\texttt{papers/gpnec/verify-all.txt}, \texttt{bench.txt},
\texttt{route-verify-paired.txt}.

Primary route table: $N=800$, $K=10$, $C=250$, $512$ packets, $80$/$120$
pre/post ticks, seed~$42$, $64$ betweenness samples.

\section{Results}
\label{sec:results}

\subsection{LBM accuracy}

\begin{table}[ht]
\centering
\caption{Metal vs CPU D2Q9 ($64\times64$, $32$ steps, $\tau=0.56$). Gate:
relative $\mathrm{L}_2\le10^{-4}$.}
\label{tab:lbm-acc}
\begin{tabular}{@{}lrrrr@{}}
\toprule
Rel.\ $\mathrm{L}_2$ & Max $|\!\Delta\!|$ & Mass rel.\ $\Delta$ & Rest drift & Gate \\
\midrule
$3.28\times10^{-7}$ & $4.77\times10^{-7}$ & $\sim10^{-6}$ & $0$ & pass \\
\bottomrule
\end{tabular}
\end{table}

\subsection{LBM throughput (scoped)}

\begin{table}[ht]
\centering
\caption{LBM step time and MLups on M4~Max (warmup~64, timed~50).
CPU columns are the kernel-faithful accuracy reference and a row-parallel
variant---not tuned CFD baselines. Metal/CPU ratios are informational only.}
\label{tab:lbm-thru}
\begin{tabular}{@{}rrrrrrr@{}}
\toprule
Size & Metal $\mu$s & Metal MLups & CPU $\mu$s & CPU-mt $\mu$s & vs CPU & $B_{\mathrm{eff}}$ est. \\
\midrule
$32^2$  & 23.4 & 43.8  & 2357  & 11280  & $101\times$ & --- \\
$64^2$  & 24.4 & 167.8 & 9633  & 47752  & $395\times$ & --- \\
$128^2$ & 34.6 & 472.9 & 38013 & 195377 & $1097\times$ & --- \\
$256^2$ & 69.9 & 937.3 & 155834& 790508 & $2229\times$ & $\sim150$~GB/s \\
\bottomrule
\end{tabular}
\end{table}

\texttt{cpu-mt} is slower than single-thread here; we include it to show that
naive multi-threading does not automatically produce a stronger baseline.

\subsection{Routing: lockstep and paired crash policies}

\begin{table}[ht]
\centering
\caption{\texttt{gpnec verify-route} checks (M4~Max, $N=800$, seed~$42$).}
\label{tab:route-verify}
\begin{tabular}{@{}lc@{}}
\toprule
Check & Outcome \\
\midrule
Embedding invariants & pass \\
Metal$\equiv$CPU hops ($8\times1024$) & $0$ / $0$ mismatches \\
Symmetric crash gate (publication) & pass \\
Sandbox crash (UI narrative; reported) & pass \\
\bottomrule
\end{tabular}
\end{table}

\begin{table}[ht]
\centering
\caption{Paired crash comparison under identical seeds
($N=800$, $K=10$, $C=250$, $512$ packets, seed~$42$).
Symmetric $=$ matched retry+respawn. Sandbox $=$ Euclidean freeze (handicapped).}
\label{tab:route-paired}
\begin{tabular}{@{}llrrr@{}}
\toprule
Policy & Metric & Post delivered & In-flight & Hyp/Euc \\
\midrule
Symmetric & Euclidean & 544 & 512 & 1.78 \\
Symmetric & Poincar\'e & 966 & 512 & 1.78 \\
Sandbox   & Euclidean & 125 & 387 & 7.65 \\
Sandbox   & Poincar\'e & 956 & 512 & 7.65 \\
\bottomrule
\end{tabular}
\end{table}

Under \textbf{symmetric} mechanics, Poincar\'e still out-delivers Euclidean
by about $1.8\times$ on this embedding and crash set---a meaningful but
far more modest gap than the sandbox's $7.7\times$, which largely reflects
disabling Euclidean retry/respawn. Pre-crash, both metrics deliver
substantial traffic (Euclidean~$888$, Poincar\'e~$3092$ cumulative under
respawn), so the post-crash comparison is not an artifact of a dead
pre-crash Euclidean mesh.

\section{Limitations and non-claims}
\label{sec:limits}

\begin{itemize}[leftmargin=*]
  \item \textbf{Routing generality.} Results are for one dual-embedding
        generator, one crash geometry, and listed seeds. We do not claim
        hyperbolic greedy routing dominates Euclidean greedy routing on
        arbitrary graphs.
  \item \textbf{Sandbox $\neq$ science claim.} Large sandbox multiples are
        expected under handicapped Euclidean policy; use symmetric numbers
        for embedding comparisons.
  \item \textbf{CPU performance.} Single-thread (and naive mt) CPU figures
        establish parity and rough scale; they are not a substitute for
        NEON/Accelerate or CUDA baselines on appropriate hardware.
  \item \textbf{Roofline.} The $\sim150$~GB/s estimate at $256^2$ is a
        back-of-envelope traffic model, not a full roofline with measured
        counters. Small lattices are likely launch/latency limited.
  \item \textbf{Cross-vendor.} No CUDA/OpenFOAM same-device study on this
        host.
  \item \textbf{Float32 hyperbolic math.} Clamping avoids NaNs near the
        disk boundary; it does not eliminate all geometric distortion from
        single-precision evaluation.
\end{itemize}

\section{Reproducibility}
\label{sec:repro}

\begin{verbatim}
swift test
swift run gpnec verify
swift run gpnec route-verify --n 800 --k 10 --crash 250 \
  --packets 512 --pre 80 --post 120 --seed 42 \
  --betweenness-samples 64 --policy both
swift run gpnec bench --backends metal,cpu,cpu-mt \
  --sizes 32,64,128,256 --steps 50 --warmup 64
\end{verbatim}

\section{Conclusion}
\label{sec:conclusion}

GPNEC's durable contributions are reproducible Metal$\equiv$CPU gates for
LBM and greedy routing, a fair symmetric crash control under which
Poincar\'e still leads on our dual embedding, and performance reporting that
separates accuracy references from competitive baselines. The asymmetric
sandbox remains a UI demonstration, not a substitute for matched-policy
evidence.

\section*{Availability}
\url{https://github.com/Digital-Defiance/gpnec}

\bibliographystyle{plainnat}
\bibliography{refs}

\end{document}
